\documentclass[12pt, a4paper]{article}

% --- PACCHETTI FONDAMENTALI ---
\usepackage[utf8]{inputenc}   % Permette di scrivere caratteri accentati direttamente da tastiera
\usepackage[T1]{fontenc}      % Migliora la qualità dei font nel PDF finale
\usepackage[italian]{babel}   % Abilita la sillabazione e i termini in italiano (es. "Sezione")

% --- GEOMETRIA E LAYOUT ---
\usepackage{geometry}         % Per una facile impostazione dei margini
\geometry{a4paper, margin=2.5cm} % Margini uniformi, adatti alla lettura
\linespread{1.2}              % Aumenta l'interlinea per una migliore leggibilità

% --- PACCHETTI UTILI ---
\usepackage{graphicx}         % Per inserire immagini (es. schemi ER)
\usepackage{amsmath}          % Per scrivere formule matematiche
\usepackage{hyperref}         % Rende i riferimenti interni (sommario, citazioni) cliccabili
\usepackage{xcolor}           % Per usare i colori (necessario per il codice)
\usepackage{listings}         % Per inserire blocchi di codice formattati
\usepackage{booktabs}         % Necessario per \toprule, \midrule e \bottomrule

\title{
    \vspace{-2cm} % Spazio negativo per alzare il titolo
    \fontsize{18pt}{22pt}\selectfont
    \textbf{Progetto - Basi di Dati}
}
\author{
    \large
    \textbf{Studenti:} \\ Giovanni Veronese \\ Sebastiano Lorenzi \\
    \vspace{2cm} \\
    \normalsize
    Prof. Massimiliano de Leoni \\
    \vspace{0.5cm}
    Corso di Laurea Informatica \\
}
\date{\textbf{Anno Accademico:} 2025-2026}

\hypersetup{
    colorlinks=true,
    linkcolor=black,      % Colore dei link interni (sommario, etc.)
    urlcolor=blue,        % Colore dei link web
    pdftitle={Relazione Progetto: Titolo del Progetto},
    pdfauthor={Nome Cognome},
    pdfsubject={Relazione per il corso di Nome Corso},
    bookmarks=true,
    bookmarksopen=true,
    bookmarksopenlevel=2
}

\definecolor{codegreen}{rgb}{0,0.6,0}
\definecolor{codegray}{rgb}{0.5,0.5,0.5}
\definecolor{codepurple}{rgb}{0.58,0,0.82}
\definecolor{backcolour}{rgb}{0.95,0.95,0.92}

\lstdefinestyle{sqlstyle}{
    backgroundcolor=\color{backcolour},
    commentstyle=\color{codegreen},
    keywordstyle=\color{magenta},
    numberstyle=\tiny\color{codegray},
    stringstyle=\color{codepurple},
    basicstyle=\ttfamily\footnotesize,
    breakatwhitespace=false,
    breaklines=true,
    captionpos=b,
    keepspaces=true,
    numbers=left,
    numbersep=5pt,
    showspaces=false,
    showstringspaces=false,
    showtabs=false,
    tabsize=2
}
\lstset{style=sqlstyle, language=SQL} % Imposta lo stile e il linguaggio di default

\begin{document}
\maketitle
\thispagestyle{empty}

\newpage

\tableofcontents
\thispagestyle{empty}

\newpage
\setcounter{page}{1} % La numerazione delle pagine riparte da 1

\section{Abstract}

In un sistema sanitario che è costantemente alla ricerca di nuovo sangue e plasma per le trasfusioni, una base di dati che garantisce una archiviazione e gestione corretta delle informazioni è essenziale, senza tralasciare sicurezza ed efficienza.

Il progetto ha la finalità di gestire un database sulle donazioni di sangue, dal momento in cui vengono create le singole sacche al momento in cui vengono utilizzate. 

Il sistema distingue i centri prelievi dagli ospedali, permettendo di registrare i prelievi effettuati dai donatori negli appositi centri di raccolta, la creazione delle sacche e la rispettiva locazione, il loro stato di disponibilità e il loro eventuale trasferimento tra ospedali. Ogni sacca può essere utilizzata in una trasfusione, collegata a una specifica richiesta di sangue e a un determinato paziente.

\section{Analisi dei requisiti}

Questa sezione riassume i requisiti a cui deve sottostare la base di dati.

Per poter donare il sangue bisogna iscriversi al database \textbf{donatore}, fornendo le seguenti informazioni:

\begin{itemize}
    \item Codice fiscale (\textbf{cf}) univoco che lo identifichi
    \item Un nominativo \textbf{nome} e \textbf{cognome}
    \item Viene assegnato un \textbf{gruppo sanguigno} [Valori possibili: 'A', 'B', 'AB', '0']
    \item L'informazione del \textbf{fattore RH} specifica il gruppo sanguigno [Valori possibili: '+', '-']
    \item Il \textbf{numero di telefono} per contattare il centro o essere contattato
\end{itemize}

Un donatore può effettuare più prelievi in momenti differenti, può donare sangue intero o plasma a piacimento e in un centro di raccolta non deciso a priori. È possibile effettuare una sola donazione alla volta e di conseguenza generare un \textbf{prelievo} identificato da:

\begin{itemize}
    \item Un \textbf{id} univoco e sequenziale
    \item Una \textbf{data} e un orario (precisione al secondo)
\end{itemize}

Ogni prelievo genera una sola sacca, e viene eseguito in un centro prelievi che rifornisce uno e un solo ospedale. Il \textbf{centro prelievi} ha i seguenti attributi:

\begin{itemize}
    \item Un \textbf{nome} che può ripetersi solo se riferisce a un diverso ospedale
    \item Un indirizzo in cui è ubicato. Contiene: \textbf{via, numero civico, CAP, città, provincia}
\end{itemize}

Il centro prelievi può eseguire più prelievi, e inviare all'ospedale di riferimento la \textbf{sacca} prodotta, il centro del progetto, che è così identificata:

\begin{itemize}
    \item Ha un univoco \textbf{id} incrementale come da norma europea
    \item Un \textbf{contenuto} [Valori possibili: 'Sangue intero', 'Plasma']. Ai fini del progetto, il sangue intero non viene diviso nelle sue emocomponenti come nella realtà.
    \item Un \textbf{gruppo sanguigno} [Valori possibili: 'A', 'B', 'AB', '0']
    \item Un \textbf{fattore RH} che accompagna il gruppo sanguigno [Valori possibili: '+', '-']
    \item Una \textbf{data di scadenza} che in questo progetto è 365 giorni per il Plasma, 40 giorni per il Sangue intero
    \item Uno \textbf{stato della sacca} per tenere traccia della disponibilità [Valori possibili: 'Disponibile', 'In uso', 'Scartata', 'In transito', 'Esaurita']
\end{itemize}

Una volta prodotta, una sacca può andare incontro a diversi destini, tracciati dallo stato: può essere immagazzinata in un ospedale (Disponibile), essere prenotata per una trasfusione (In uso), scadere o venire arbitrariamente tolta dalle scorte per motivi medici (Scartata), essere oggetto a trasferimento in un altro ospedale (In transito), o essere utilizzata per una trasfusione (Esaurita).
Di default, viene immediatamente conservata con stato 'Disponibile' nell'ospedale a cui il centro prelievi fa riferimento. Ogni \textbf{ospedale} ha i seguenti attributi:

\begin{itemize}
    \item Un \textbf{id} univoco
    \item Un \textbf{nome} che non ha scopi logici
    \item Un indirizzo per i trasferimenti. Contiene: \textbf{via, numero civico, CAP, città, provincia}
\end{itemize}

Ogni ospedale ha un proprio deposito per immagazzinare le sacche sovradimensionato rispetto al numero di sacche (non necessario controllare il numero massimo di stoccaggio). Può ricevere o spedire le sacche, ricoverare pazienti e richiedere il sangue. La \textbf{richiesta di sangue} è identificata da:

\begin{itemize}
    \item 
\end{itemize}


\section{Progettazione concettuale}
\section{Progettazione logica}
\section{Implementazione PostgreSQL e query}
\section{Codice C}














\section{Introduzione}
\label{sec:introduzione}
Introduzione al progetto, idea ecc...

\begin{table}[h]
\centering
\begin{tabular}{llcl}
\toprule
\textbf{Attributo} & \textbf{Chiave} & \textbf{Tipo} & \textbf{Vincoli e Dominio} \\
\midrule
cf\_donatore & PK & CHAR(16) &  Regex alfanumerico \\
nome            && VARCHAR(128) & NOT NULL \\
cognome          & & VARCHAR(128) & NOT NULL \\
gruppo\_sanguigno & & ENUM         & Valori: 'A', 'B', 'AB', '0' \\
fattore\_rh       & & CHAR(1)      & IN ('+', '-') \\
numero\_di\_telefono & & VARCHAR(15)  & Formato E.164 \\
\bottomrule
\end{tabular}
\caption{Dizionario dei dati: Entità DONATORE}
\end{table}

\subsection{Esempio di Query SQL}
Esempio di codice SQL

\begin{lstlisting}[caption={Esempio di riquadro in SQL}, label={}]
-- Seleziona il nome del paziente e conta le sue visite
SELECT
    p.Nome,
    p.Cognome,
    COUNT(v.ID_Visita) AS NumeroVisite
FROM Paziente AS p
JOIN Visita AS v ON p.ID_Paziente = v.ID_Paziente
WHERE v.DataVisita >= '2023-01-01'
GROUP BY p.ID_Paziente, p.Nome, p.Cognome
ORDER BY NumeroVisite DESC;
\end{lstlisting}

\section{Conclusioni}
\label{sec:conclusioni}
Testo conclusioni
\end{document}
